把叶子数补到二次幂 $s$，建立两棵迭代线段树：出树从父结点向子结点连 $0$ 权边，用于点到区间；入树从子结点向父结点连 $0$ 权边，用于区间到点。出树叶子就是原图顶点，并向对应的入树叶子连 $0$ 权边。

\textbf{只在接口层使用原图编号。}原图共有 $n$ 个点，对外统一编号为 $1,2,\ldots,n$。不要直接用原图编号访问 \texttt{g}；原图点 $u$ 在扩展图中的真实编号是 \texttt{id(u)}。线段树结点和中转点由结构体内部管理。

接口约定：

\begin{itemize}
  \item \texttt{edge(u,v,w)}：添加普通边 $u\to v$；
  \item \texttt{point\_range(u,l,r,w)}：添加 $u\to[l,r]$ 中每个点的边；
  \item \texttt{range\_point(l,r,v,w)}：添加 $[l,r]$ 中每个点到 $v$ 的边；
  \item \texttt{range\_range(l1,r1,l2,r2,w)}：添加第一个区间到第二个区间的边。
\end{itemize}

一次区间连边只覆盖迭代线段树上的 $O(\log n)$ 个结点。区间到区间通过一个中转点连接两侧，避免把覆盖结点两两连边，因此仍只新增 $O(\log n)$ 条边。

扩展图保存在 \texttt{g} 中。跑图算法时，源点写为 \texttt{id(s)}，最后读取原图点 $u$ 的结果时使用 \texttt{d[id(u)]}。线段树内部新增的边权均为 $0$。
